\documentclass[11pt,letterpaper]{letter}
\usepackage[T1]{fontenc}
\usepackage[utf8]{inputenc}
\usepackage[margin=1in]{geometry}
\usepackage{hyperref}
\usepackage{parskip}

\signature{Ihor Kendiukhov}
\address{Ihor Kendiukhov \\ Department of Computer Science \\ University of Tuebingen \\ Tuebingen, Germany \\ ihor.kendiukhov@student.uni-tuebingen.de}

\begin{document}

\begin{letter}{The Editors \\ PLOS ONE}

\opening{Dear Editors,}

I am pleased to submit my manuscript entitled \textbf{``MI-Workbench: an autonomous multi-agent platform for mechanistic interpretability of biological foundation models''} for consideration for publication in PLOS ONE as a research article.

Biological foundation models---single-cell transformers such as Geneformer, scGPT, and scBERT---are increasingly used to make claims about gene regulation and cellular biology, yet the mechanistic interpretability (MI) studies that probe what these models actually learn remain manual, one-off, and hard to reproduce. Such investigations demand repeated cycles of hypothesis formation, controlled experimentation, statistical evaluation, and adversarial critique. This manuscript presents MI-Workbench, an open-source platform that automates these cycles through orchestrated multi-agent research loops.

The key methodological contributions of the work are:

\begin{itemize}
  \item A graph-based workflow domain-specific language (DSL) for defining configurable multi-agent research loops with conditional execution and explicit stop conditions.
  \item A consensus reviewer mechanism that runs independent critic agents in parallel and merges their findings through Jaccard-similarity deduplication and agreement-based severity escalation.
  \item A structured, artifact-aware feedback mechanism that maps reviewer critiques to specific artifacts and sections, enabling targeted revision rather than wholesale rewriting.
  \item A provider-agnostic adapter layer abstracting over multiple large-language-model backends, and an end-to-end reproducibility-packaging pipeline that bundles artifacts, prompts, environment metadata, and provenance.
\end{itemize}

We validate the platform in a case study asking whether attention heads in Geneformer encode causal gene regulatory relationships or merely reflect co-expression, in which the system autonomously identified well-documented confounders (co-expression bias, rank-value encoding artifacts, and attention symmetry) and iteratively refined its analysis without human intervention. The system is implemented in over 10,000 lines of Python and is accompanied by a suite of 360 automated tests.

I believe this work is well suited to PLOS ONE. It is a technically rigorous, methodologically complete contribution whose conclusions are supported by the results presented, and it places a strong emphasis on reproducibility: the complete source code, prompt templates, loop definitions, and test suite are openly available, and the platform generates self-contained reproducibility packages for every run. These features align directly with PLOS ONE's publication criteria and its commitment to open, reproducible science.

I confirm that this manuscript is original, has not been published previously, and is not under consideration for publication elsewhere. There are no prior or concurrent submissions derived from this work. As the sole author, I have read and approved the manuscript, and I have no competing interests to declare. I received no specific funding for this work. All data and software underlying the findings are openly available, as described in the manuscript's Supporting Information.

Thank you for considering this submission. I look forward to the reviewers' feedback and would be happy to provide any additional information required.

\closing{Sincerely,}

\end{letter}
\end{document}
